\documentclass[11pt]{article}
\usepackage[margin=1in]{geometry}
\usepackage{amsmath,amssymb,amsthm}
\usepackage{graphicx}
\usepackage[hidelinks]{hyperref}

\newtheorem{theorem}{Theorem}
\newtheorem{lemma}{Lemma}
\newtheorem{remark}{Remark}

\newcommand{\R}{\mathbb R}
\newcommand{\C}{\mathbb C}
\newcommand{\K}{\mathbb K}
\newcommand{\tr}{\operatorname{tr}}
\newcommand{\OFP}{\operatorname{OFP}}

\title{Dimension Two Has No Spurious Local Minima for the Operator-Valued Frame Potential}
\author{fa\_banach\_001, agent\_lane\_16}
\date{2026-07-03}

\begin{document}
\maketitle

\section*{Source Conjecture}

This packet concerns Conjecture 5.1 of Needham--Shonkwiler, \emph{Fusion Frame Homotopy and Tightening Fusion Frames by Gradient Descent}, arXiv:2208.11045. In the authors' notation, for fixed spectra \(\mathbf r\), the conjecture states that all local minima of the obvious operator-valued frame potential
\[
        \OFP:\mathcal{OF}^{\K^d,\mathbf k}(\mathbf r)\to\R
\]
are global minima.

\begin{center}
\includegraphics[width=0.98\textwidth]{figures/open_problem_crop.png}
\end{center}

\section*{Result}

We prove the conjecture for \(d=2\), over both \(\K=\R\) and \(\K=\C\), with arbitrary finite fixed spectral data. This is a partial result: it does not address the higher-dimensional conjecture.

\section*{Proof Intuition}

In dimension two, a Hermitian operator with fixed eigenvalues is just a scalar matrix plus a fixed-radius traceless part. The traceless part is parametrized by a unit circle in the real case and by a unit sphere in the complex case. Therefore the variable part of \(\OFP\) is the squared length of a weighted vector sum
\[
        \left|\sum_i c_i x_i\right|^2,
        \qquad x_i\in S^q,\quad c_i\ge 0.
\]
This elementary energy has no spurious local minima. If the weighted sum is zero, the point is globally minimizing. If it is nonzero, first variation forces every \(x_i\) to be parallel or antiparallel to the sum. A collinear critical point with two or more vectors pointing along the sum can be lowered by rotating two of them in opposite tangent directions. Hence a nonzero local minimum can only occur when one weight dominates all the others; by the triangle inequality that configuration is already globally minimizing.

\section*{Weighted-Sphere Lemma}

\begin{lemma}\label{lem:weighted-sphere}
Let \(q\ge 1\), let \(c_1,\ldots,c_N\ge 0\), and define
\[
        F:(S^q)^N\to\R,\qquad
        F(x_1,\ldots,x_N)=\left|\sum_{i=1}^N c_i x_i\right|^2 .
\]
Then every local minimum of \(F\) is a global minimum.
\end{lemma}

\begin{proof}
Discarding the indices with \(c_i=0\) changes neither local nor global minimality, so assume \(c_i>0\). Let \(x=(x_i)\) be a local minimum and set
\[
        R=\sum_i c_i x_i.
\]
If \(R=0\), then \(F(x)=0\), so \(x\) is a global minimum.

Assume \(R\ne 0\), and write \(u=R/|R|\). The first variation of \(F\) in the tangent space \(T_{x_i}S^q\) is
\[
        v_i\mapsto 2c_i\langle R,v_i\rangle .
\]
At a local minimum this vanishes for every \(v_i\perp x_i\). Hence \(R\) is parallel to each \(x_i\), so \(x_i=\varepsilon_i u\) with \(\varepsilon_i\in\{\pm1\}\). Put
\[
        L=\sum_i \varepsilon_i c_i=|R|>0.
\]

We claim that at most one sign \(\varepsilon_i\) can be \(+1\). Suppose instead that \(\varepsilon_p=\varepsilon_q=+1\) for two distinct indices \(p,q\). Choose a unit vector \(v\perp u\), possible since \(q\ge 1\), and choose nonzero real numbers \(a_p,a_q\) with \(c_p a_p+c_q a_q=0\). Consider the curve that leaves all \(x_i\) fixed except
\[
        x_p(t)=u\cos(a_p t)+v\sin(a_p t),\qquad
        x_q(t)=u\cos(a_q t)+v\sin(a_q t).
\]
The first-order tangent change of \(R\) is zero, while the second-order change in the \(u\)-direction is
\[
        -\frac{t^2}{2}\bigl(c_p a_p^2+c_q a_q^2\bigr)u+O(t^3).
\]
Consequently
\[
        F(x(t))=L^2-L\bigl(c_p a_p^2+c_q a_q^2\bigr)t^2+O(t^3)<L^2
\]
for all sufficiently small nonzero \(t\), contradicting local minimality.

Thus exactly one index, say \(p\), has \(\varepsilon_p=+1\), and
\[
        L=c_p-\sum_{i\ne p}c_i>0.
\]
For any \(y_i\in S^q\), the triangle inequality gives
\[
        \left|\sum_i c_i y_i\right|
        \ge c_p-\sum_{i\ne p}c_i=L.
\]
Therefore \(x\) attains the global minimum value. This proves the lemma.
\end{proof}

\section*{Dimension-Two Theorem}

\begin{theorem}\label{thm:d2}
Let \(\K=\R\) or \(\C\). In Hilbert dimension \(d=2\), for arbitrary finite fixed spectral data \(\mathbf r\), every local minimum of the generalized operator-valued frame potential on \(\mathcal{OF}^{\K^2,\mathbf k}(\mathbf r)\) is a global minimum.
\end{theorem}

\begin{proof}
It is enough to work with the positive semidefinite operators
\[
        P_i=A_i^*A_i
\]
with prescribed spectra; the unitary ambiguity in the square roots \(A_i\) does not change \(\OFP\). Thus
\[
        \OFP(P_1,\ldots,P_N)=\left\|\sum_i P_i\right\|_{\mathrm{HS}}^2
        =\tr\left(\left(\sum_i P_i\right)^2\right).
\]

Fix \(i\). Let the two eigenvalues of \(P_i\) be \(\alpha_i\ge \beta_i\ge 0\), where \(\beta_i=0\) is allowed for rank-one data. Write
\[
        m_i=\frac{\alpha_i+\beta_i}{2},
        \qquad
        c_i=\frac{\alpha_i-\beta_i}{2}.
\]
Then
\[
        P_i=m_i I_2+c_i X_i,
\]
where \(X_i\) is self-adjoint, traceless, and has eigenvalues \(1,-1\), unless \(c_i=0\), in which case \(P_i=m_iI_2\) is constant and can be ignored.

The set of such \(X_i\) is naturally a unit sphere: for \(\K=\C\), it is the Bloch sphere in the three-dimensional real space of traceless Hermitian \(2\times2\) matrices; for \(\K=\R\), it is a unit circle in the two-dimensional real space of traceless symmetric \(2\times2\) matrices. In either case, after choosing an orthonormal basis of the traceless space, write \(X_i\leftrightarrow x_i\in S^q\), where \(q=2\) in the complex case and \(q=1\) in the real case.

Now
\[
        \sum_i P_i=\left(\sum_i m_i\right)I_2+\sum_i c_iX_i.
\]
The scalar and traceless summands are Hilbert--Schmidt orthogonal, so
\[
        \OFP(P_1,\ldots,P_N)
        =
        \text{constant}
        +
        C\left|\sum_i c_i x_i\right|^2
\]
for a fixed positive normalization constant \(C\) depending only on the chosen inner-product convention. Multiplication by a positive constant and addition of a constant preserve local and global minimizers. Therefore the local minima of \(\OFP\) are exactly the local minima of the weighted-sphere energy in Lemma~\ref{lem:weighted-sphere}. By that lemma, they are global minima.
\end{proof}

\section*{Verification Notes}

The proof is analytic and self-contained. A short exploratory Hessian search over diagonal critical points in \(d=3\) found no counterexample, but no computational claim is used here and no higher-dimensional conclusion is asserted. The theorem depends only on the elementary two-dimensional orbit parametrization and Lemma~\ref{lem:weighted-sphere}.

\section*{Novelty Check and Limitations}

Local run indexes were searched for the arXiv id and close phrases around operator-valued frame potentials, local versus global minima, and the fusion-frame Benedetto--Fickus theorem. Web searches were also run for close variants of those phrases and returned the source paper but no independent proof of this \(d=2\) case. Novelty confidence is moderate pending expert review.

This packet proves only the \(d=2\) case of Conjecture 5.1. The conjecture for \(d\ge 3\), including the full operator-valued fixed-spectrum setting, remains open.

\section*{References}

\begin{thebibliography}{9}
\bibitem{NeedhamShonkwiler}
T. Needham and C. Shonkwiler,
\emph{Fusion Frame Homotopy and Tightening Fusion Frames by Gradient Descent},
arXiv:2208.11045, 2022.
\end{thebibliography}

\end{document}
